\section{JOINT-EVIDENCE MODEL AND COMPARISON FRAMEWORK}
\label{sec:foundations_comparison}
The value of a cross-platform O-ISAC review does not come from placing studies
under platform labels alone. Its central analytical task is to determine when
a reported communication result and a sensing result describe a joint
engineering relation, which conditions bound that relation, and which
conclusion the evidence permits. We therefore use a relation-centered unit of
synthesis. This section defines that unit, specifies the measurement and
evidence conditions carried with it, and converts those conditions into
explicit comparison permissions. Section~\ref{sec:review_methods} then explains
how the framework is populated from the review corpus.

\subsection{What Counts as a Joint-Performance Relation?}

Two performance values do not become joint evidence merely because they
appear in the same paper or originate from an integrated architecture. The
minimum native relation requires a shared or varied design variable,
constraint, or impairment to be linked to both the communication and sensing
outcomes under an adequate measurement and evidence contract. A complete
record that can be tested for transfer additionally resolves the physical
system and the coupling mechanism. For relation record $r$, we express this
full unit as
\[
J_r=\langle P_r,G_r,X_r,C_r,S_r,M_r,E_r\rangle .
\]
Here, $P_r$ identifies the physical context and system boundary; $G_r$
locates what the communication and sensing functions share and where they
share it; $X_r$ is the shared or varied design variable, constraint, or
impairment linked to both functions; $C_r$ and $S_r$ are the reported
communication and sensing outcomes; and $M_r$ and $E_r$ preserve the
measurement contract and evidence envelope. The tuple is not a score, and the
presence of a native link does not by itself establish transferability. It is a
compact test of whether a claimed relation can be interpreted and assessed
without changing the meaning of either result.

The physical context is set by the path that reaches the target, not by the
presence of optical hardware alone. The observable modality is therefore kept
separate from the technology used to generate, transport, or process the
signal. A fiber may itself be the sensing medium or may transport a signal to
a downstream target-facing segment. In a photonic-terahertz system, optical
components may generate or process the carrier while the target interaction
occurs over a radio-frequency path. The corresponding propagation, conversion,
illumination, and geometry boundaries are retained for all platform families
before comparison in Section~\ref{sec:optical_modality_families}
\cite{OISAC_SCR00220,OISAC_SCR00940,OISAC_SCR00196,OISAC_SCR00083,
OISAC_SCR00366}.

The coupling term $G_r$ records where the functions meet; several coupling
loci may coexist. The factor $X_r$ is kept separate from that shared object: it
is the variable, constraint, or impairment whose setting or exposure is
related to both outcomes. This distinction separates \emph{structural
integration} from \emph{evidential jointness}. Shared hardware or a common
service establishes an architecture relation, whereas an adequate
$X_r$--$C_r$--$S_r$ trace establishes the minimum native performance relation.
Resolving $P_r$ and $G_r$ is additionally required before the complete record
can be tested for transfer beyond its native context.
Figure~\ref{fig:joint_evidence_anatomy} shows this anatomy.

\begin{figure*}[!t]
\centering
\footnotesize
\setlength{\tabcolsep}{2.5pt}
\renewcommand{\arraystretch}{1.08}
\begin{tabular}{@{}c@{\hspace{5pt}}c@{\hspace{5pt}}c@{\hspace{5pt}}c@{\hspace{5pt}}c@{}}
\fbox{\parbox[c][13mm][c]{0.275\textwidth}{\centering
\textbf{Minimum native relation}\\
$X_r$ linked to $C_r$ and $S_r$ under adequate $M_r/E_r$}}
& $+$ &
\fbox{\parbox[c][13mm][c]{0.225\textwidth}{\centering
\textbf{Resolved context}\\
physical boundary $P_r$ and coupling $G_r$}}
& $\longrightarrow$ &
\fbox{\parbox[c][13mm][c]{0.25\textwidth}{\centering
\textbf{Complete record $J_r$}\\
eligible for the transfer test}}
\\
\end{tabular}
\caption{Native-relation and complete-record requirements. The minimum native
relation links a shared or varied factor to both outcomes under an adequate
measurement contract and evidence envelope. Resolving the physical and
coupling context completes $J_r$ and permits assessment of transfer beyond the
source setting; transfer permission still depends on adequate, aligned $M/E$
conditions and independent evidence.}
\label{fig:joint_evidence_anatomy}
\end{figure*}

A native link is most directly established within one study because its two
outcomes can be followed through the same model, experiment, dataset, or
operating sweep. Transfer across studies is a separate operation. It is
allowed only after the meanings and conditions attached to the native
relations have been aligned; co-occurrence at the topic or platform level is
not sufficient.

\subsection{Measurement and Evidence Conditions}

The link in $J_r$ remains interpretable only while the measurement and
evidence contracts remain attached to it. We organize them as
\[
M_r=\langle M_C,M_S,A_M\rangle,
\qquad
E_r=\langle E_C,E_S,T_r,B_r\rangle .
\]
$M_C$ and $M_S$ retain the definitions used for the communication and
sensing quantities, respectively. Each includes the measurement plane, unit,
statistic or accounting state, task and target, operating point, channel,
geometry, processing conditions, baseline, and ground truth. $A_M$ records
whether the condition sets are the same, alignable, materially different, or
not yet resolved for comparison. This
alignment term prevents a familiar metric name from substituting for an
actual measurement match.

Several distinctions are particularly consequential in optical systems. OSNR
and electrical SNR refer to different observation planes; launch and received
power refer to different positions in the signal chain. Gross and net rate,
or error rates before and after correction, use different accounting states.
For sensing, resolution, estimation error, precision, detection probability,
and an analytical bound answer different questions. A numerical comparison is
therefore formed from matched definitions and references, not from matched
abbreviations.

$E_C$ and $E_S$ identify how each outcome was produced, the validation setting,
and the exact source locator at which it can be verified. Analytical,
simulation, laboratory, offline-dataset, and field evidence are not treated as
interchangeable; each supports inference within its own validation setting.
$T_r$ preserves the configuration or timing trace that connects the outcomes
to $X_r$, such as a shared sweep, timestamped observation, configuration
identifier, or optimization run. $B_r$ states the permitted inference boundary:
the native setting or aligned transfer domain within which the relation may be
used. A validation label by itself does not create a paired relation: the link
must still be recoverable through $T_r$.

These conditions are deliberately separated from a study-level validation
label. A nontransferable result may still be decisive for its source system,
and field validation does not repair a mismatch in metric meaning. The
framework instead asks a narrower question: what is the strongest conclusion
that follows while the original system, measurement, and evidence meanings
are preserved? Figure~\ref{fig:claim_permission_pipeline} summarizes that
decision path.

\begin{figure*}[!t]
\centering
\footnotesize
\setlength{\tabcolsep}{2.5pt}
\begin{tabular}{@{}c@{\hspace{2pt}}c@{\hspace{2pt}}c@{\hspace{2pt}}c@{\hspace{2pt}}c@{\hspace{2pt}}c@{\hspace{2pt}}c@{\hspace{2pt}}c@{\hspace{2pt}}c@{}}
\fbox{\parbox[c][12mm][c]{0.14\textwidth}{\centering
\textbf{1. Locate}\\
physical regime and coupling}}
& $\longrightarrow$ &
\fbox{\parbox[c][12mm][c]{0.14\textwidth}{\centering
\textbf{2. Link}\\
trace $X_r$ to both $C_r$ and $S_r$}}
& $\longrightarrow$ &
\fbox{\parbox[c][12mm][c]{0.145\textwidth}{\centering
\textbf{3. Align}\\
measurement contracts and condition sets}}
& $\longrightarrow$ &
\fbox{\parbox[c][12mm][c]{0.145\textwidth}{\centering
\textbf{4. Bound}\\
provenance, trace, and inference scope}}
& $\longrightarrow$ &
\fbox{\parbox[c][12mm][c]{0.145\textwidth}{\centering
\textbf{5. State}\\
the permitted claim class}}\\
\end{tabular}
\caption{Claim-permission path used in the review. Each stage preserves the
meaning established by the previous stage. An unmet gate does not erase the
source result; it routes the record to the strongest permitted class in
Table~\ref{tab:comparison_record}.}
\label{fig:claim_permission_pipeline}
\end{figure*}

\subsection{From a Relation Record to a Permitted Claim}

The framework converts evidence description into an explicit synthesis
decision. The weakest material condition bounds the claim; a
strong value on one metric cannot compensate for a missing link or an
incompatible definition elsewhere in the relation. Table~\ref{tab:comparison_record}
states the resulting permission matrix. The matrix separates source-native
performance, within-study relations, and cross-study transfer so that the
reader can see both what the evidence supports and where that support ends.

\begin{table*}[!t]
\caption{Evidence conditions and permitted claim classes in the O-ISAC synthesis}
\label{tab:comparison_record}
\centering
\footnotesize
\setlength{\tabcolsep}{3.2pt}
\renewcommand{\arraystretch}{1.12}
\begin{tabularx}{\textwidth}{@{}>{\raggedright\arraybackslash}p{0.22\textwidth}
>{\raggedright\arraybackslash}p{0.19\textwidth}
>{\raggedright\arraybackslash}p{0.25\textwidth}
>{\raggedright\arraybackslash}X@{}}
\toprule
Evidence condition & Permitted claim & Claim boundary & Downstream use \\
\midrule
An integration mechanism or shared element is documented, but no traceable
$X_r$--$C_r$--$S_r$ link is established. &
Architecture or feasibility statement. &
Shows how the functions can coexist or share resources; it does not establish
a joint-performance relation. &
Architecture synthesis in Section~\ref{sec:optical_modality_families}. \\

A communication or sensing result is source-traceable and its native
conditions are recoverable. &
Source-native performance statement. &
Applies to the reported task, system, and operating region; no effect on the
other function is inferred. &
Metric landscape in Section~\ref{sec:metrics}. \\

One study links $X_r$ to both outcomes through a common trace, $A_M$ confirms
the same or alignable condition set, and the $M/E$ contract is adequate. &
Within-study relation or tradeoff. &
The source-supported direction and any source-reported magnitude apply only to
the tested model, experiment, dataset, and operating region. &
Native tradeoff synthesis in Section~\ref{sec:tradeoffs}. \\

At least two independent studies expose compatible native relations and align
physical regime, metric definitions, planes, units, tasks, targets, operating
conditions, baselines, validation setting, and relevant uncertainty; remaining
differences are bounded. &
Bounded cross-study numerical comparison. &
Transfer is limited to the shared comparison domain; it is not a universal
effect, unconditional pooled estimate, platform ranking, or leaderboard. &
Cross-platform synthesis in Section~\ref{sec:tradeoffs}. \\

A recurrent, source-supported coupling interpretation connects both functions,
but the numerical contracts do not support magnitude transfer. &
Mechanism-based or directional statement. &
States a source-supported direction or recurring coupling interpretation
without assigning a shared numerical effect or unsupported causality. &
Mechanism synthesis across Sections~\ref{sec:optical_modality_families}
and~\ref{sec:tradeoffs}. \\

A recurring, material omission prevents a decision-relevant relation from
being tested or transferred. &
Actionable literature-gap statement. &
The gap names the missing measurement, baseline, trace, or validation link and
the comparison it is designed to test or enable. &
Benchmark priorities in Section~\ref{sec:validation_reproducibility} and the
roadmap in Section~\ref{sec:discussion_roadmap}. \\

No supported relation route is available, or the available routes remain in
material conflict. &
Audit retention only. &
The record remains in the evidence base for transparency and future updating;
no reader-facing relational conclusion is drawn from it. &
Corpus accounting in Section~\ref{sec:review_methods}. \\
\bottomrule
\end{tabularx}
\end{table*}

When either the independence test or the alignment test for numerical transfer
is not satisfied, the decision is explicit non-pooling. The most informative
supported non-cross-study use is retained, but the review does not produce a
pooled effect, forced unit normalization, platform ranking, or leaderboard.
This is a comparison decision, not a judgment that the underlying studies are
uninformative. The framework likewise does not manufacture a positive relation
merely to fill a synthesis category.

This permission structure supplies the handoff to the remainder of the review.
Section~\ref{sec:review_methods} constructs and traces the records;
Section~\ref{sec:optical_modality_families} explains their physical and
architectural mechanisms; Section~\ref{sec:tradeoffs} evaluates the admitted
performance relations; Section~\ref{sec:validation_reproducibility} assesses
their reconstructability and benchmark readiness;
Section~\ref{sec:technologies_applications_6g} translates supported mechanisms
to application and 6G requirements; and Section~\ref{sec:discussion_roadmap}
converts recurring comparison limits into
benchmark and research priorities. The result is a framework that tells an
O-ISAC researcher not only where a study belongs, but also which findings can
be compared, why that comparison is defensible, and what evidence would enable
the next stronger conclusion.

\FloatBarrier
